\chapter{综合讨论与总结}
\label{chap:discussion}

\section{同类技术在四个任务中的作用}
\begin{table}[H]
\centering\small
\caption{四项任务的核心方案与榜单口径}\label{tab:overview}
\begin{tabularx}{\linewidth}{@{}P{29mm}YP{32mm}P{22mm}P{22mm}@{}}
\toprule
任务 & 选定路径 & 指标 & Public & Private \\
\midrule
植物分类 & 修复裁剪，ConvNeXt / EfficientNet 集成 & Kaggle 得分 & 0.9735 & 0.9735 \\
表情分类 & ConvNeXt 多骨干集成与概率 bagging 并列 & 准确率 & 0.7470 & 0.7470 \\
中心凹定位 & 血管热图，两种子与 VascX 融合 & 坐标 RMSE（px） & 9.90927 & 19.54239 \\
血管分割 & FR-UNet 权重与 DTI、FOV 过滤 & Dice & 0.81766 & 0.82534 \\
\bottomrule
\end{tabularx}
\end{table}

四项任务均利用图像模型，但监督信号位于不同尺度。分类的损失作用于图像级
概率，中心凹使用空间热图和坐标误差，分割则要求逐像素保留狭窄结构。
这直接影响输入尺寸、增强和推理的选择。植物的俯拍姿态支持旋转，
人脸表情依赖局部五官结构，眼底血管图和中心凹坐标需要与原图几何一致。
植物裁剪问题说明，几何处理应从像素数据和坐标转换两方面核查，
同一训练曲线可能无法揭示系统性输入错误。

迁移学习在小样本任务中提供可用的视觉表征。植物和表情任务重新训练分类头，
中心凹任务使用预训练编码器并学习热图解码器。分割的最终方案直接采用上游
同域 FR-UNet 权重，自训 U-Net 则作为不同输入表示和权重来源的对照。
这四项任务使用了已有骨干、损失和集成机制。课程工作中的技术贡献主要是
针对任务的适配、可定位的工程缺陷修复以及用实际对照支持模型选择。

\section{实验对照与可迁移性}
表~\ref{tab:overview} 中的数值保留原任务单位，各项目的比较在同一赛题
内部进行。植物修复可以沿着错误坐标、暗区图像、重建缓存和重新训练逐级检查。
表情的分辨率对照把计算开销和精度放在同一张表中，显示增加尺寸的收益
需要结合训练预算。中心凹多折实验中，一个难例可主导该折平方误差，
因而 OOF 均值之外还应查看分折分布。血管分割则对阈值、视野与度量区域
敏感。把不同评价范围的 Dice 直接相减会掩盖模型和后处理之间的差异。

集成在四项任务中的收益并不一致。中心凹的多个来源输出能互相校正，
而表情 logit 融合的本地增益没有在榜单复现。原生 U-Net 与 FR 的融合
低于冻结 FR，说明引入另一模型会同时引入其偏差。当前实验均使用课程
固定划分，因果归因需要在控制骨干、数据处理和选择规则时再做检验。
本报告的效果陈述限于可追溯的现有对照。通过源码、配置和提交编号
可以核查完整的实验链，并对缺失的历史生成细节保留记录。

\section{复现边界}
\begin{table}[H]
\centering\small
\caption{主要结论所依赖的比较条件}
\begin{tabularx}{\linewidth}{@{}P{23mm}YY@{}}
\toprule
任务 & 可核查的对照条件 & 外推受限的原因 \\
\midrule
幼苗 & 同一骨干和分辨率的修复前后记录 & 测试图伪标签与固定测试池绑定，额外同源标注需要去重 \\
表情 & 同折、同骨干的分辨率与耗时对照 & 硬件、批大小及不同来源人脸图的分布会改变权衡 \\
定位 & 同一组折外坐标和两侧榜单均有记录 & 少量难例、重复图分组和融合阈值使各折精度波动 \\
分割 & 同一权重的多个后处理设置 & 图像缩放、概率阈值及掩膜范围共同决定假阳性和漏检 \\
\bottomrule
\end{tabularx}
\end{table}

归档的模型训练日志、折划分和榜单数据支持逐项核对本报告。完整重训需要
课程数据访问、上游权重和适当的运行资源。公开源码保留关键入口与配置。
记录中缺少的历史命令、权重来源细节和部分超参数在证据表中标记。正式报告
中的样例均由本地授权访问的课程文件和已保存输出生成。写作与证据整理使用
Scientific Agent Skills 提供的流程辅助\citep{kassis2026skills}。
